\documentclass[12pt,a4paper]{article}
\usepackage[utf8]{inputenc}
\usepackage[italian]{babel}
\usepackage{amsmath, amssymb, amsfonts}
\usepackage{graphicx}
\usepackage{pgfplots}
\pgfplotsset{compat=1.18}
\usepackage{geometry}
\geometry{margin=2.5cm}
\usepackage{caption}
\usepackage{hyperref}

\title{\centering Studio matematico sulla coscienza umana \\ Legge 4: Legge Generale della Coscienza Quadristato}
\author{Autore: Giuseppe Carlino}
\date{\today}

\begin{document}

\maketitle

\section{Premessa}

La \textbf{Legge Generale della Coscienza Quadristato} unifica i quattro stati fondamentali della coscienza:

\[
\mathbf{X} = 
\begin{pmatrix}
t_c \\ k \\ s \\ w
\end{pmatrix},
\]

dove:

\begin{itemize}
    \item $t_c$ è il tempo percepito (Legge 1),
    \item $k$ è la conoscenza accumulata (Legge 0),
    \item $s$ è la scelta effettuata (Legge 2),
    \item $w$ è la volontà risultante (Legge 3).
\end{itemize}

La legge generale formalizza l’evoluzione della coscienza come un sistema dinamico quadridimensionale, soggetto a perturbazioni e fluttuazioni interne ed esterne, che integrano conoscenza, tempo, scelta e volontà.

\section{Formalizzazione matematica}

Si definisce un \textbf{equilibrio assoluto} della coscienza:

\[
O = 
\begin{pmatrix}
t_e \\ k_e \\ s_e \\ w_e(F)
\end{pmatrix},
\]

e un \textbf{coefficiente di squilibrio}:

\[
\kappa = \sqrt{(t_c - t_e)^2 + (k - k_e)^2 + (s - s_e)^2 + (w - w_e)^2}.
\]

L’evoluzione dinamica dei quattro stati è descritta dal sistema differenziale:

\begin{equation}
\frac{d\mathbf{X}}{d\tau} = - \Lambda \cdot (\mathbf{X} - O) + \mathbf{\eta}(\tau),
\end{equation}

dove:

\[
\Lambda = 
\begin{pmatrix}
\lambda & 0 & 0 & 0 \\
0 & \mu & 0 & 0 \\
0 & 0 & \nu & 0 \\
0 & 0 & 0 & \xi
\end{pmatrix}, 
\quad
\mathbf{\eta}(\tau) =
\begin{pmatrix}
\eta_t \\ \eta_k \\ \eta_s \\ \eta_w
\end{pmatrix}.
\]

\begin{itemize}
    \item $\lambda, \mu, \nu, \xi > 0$ sono coefficienti di richiamo verso l’equilibrio relativo dei singoli stati.
    \item $O$ rappresenta l’equilibrio assoluto della coscienza.
    \item $\mathbf{\eta}(\tau)$ rappresenta perturbazioni interne o esterne.
    \item $w_e(F)$ evidenzia la dipendenza della volontà dai fattori multipli $F$, quindi la complessità del sistema.
\end{itemize}

\section{Diagramma concettuale quadridimensionale}

\begin{center}
\begin{tikzpicture}
\begin{axis}[
    view={70}{25},
    width=14cm, height=10cm,
    axis lines=center,
    grid=major,
    tick label style={font=\small},
    xmin=0, xmax=6,
    ymin=0, ymax=4,
    zmin=0, zmax=5,
    enlargelimits=0.2,
    axis line style={->},
    ticks=none,
    xlabel={$t_c$}, ylabel={$k$}, zlabel={$w$}
]

% Traiettoria concettuale della coscienza
\addplot3[smooth, thick, blue] coordinates {
(0,0.5,0.2) (0.5,0.8,0.5) (1,1.0,1.0) (1.5,1.2,1.5) (2,1.5,2.0)
(2.5,1.7,2.3) (3,2.0,2.5) (3.5,2.3,3.0) (4,2.5,3.2) (4.5,2.7,3.5) (5,3.0,4)
};

% Punto equilibrio assoluto O
\addplot3[only marks, red, mark=*] coordinates {(3,2,2.5)};
\node at (axis cs:2.2,2.5,4.3) [anchor=west] {\small $O=(t_e,k_e,s_e,w_e)$};

% Nuvola concettuale: punti statici intorno a O
\addplot3[only marks, mark=*, mark size=0.7pt, gray, opacity=0.3] coordinates {
(3.1,2.1,2.6) (2.9,1.9,2.4) (3.05,2.0,2.5) (2.95,2.05,2.45) (3.0,2.0,2.55)
};

% Nota interpretativa
\node at (axis cs:2.5,1.5,5.3) [anchor=west] {\scriptsize Fluttuazioni e traiettoria concettuale attorno a $O$};

\end{axis}
\end{tikzpicture}
\end{center}

\caption*{Lo spazio quadridimensionale (rappresentato tridimensionalmente) mostra la coscienza umana integrando tempo $t_c$, conoscenza $k$, scelta $s$ e volontà $w$. La traiettoria blu rappresenta l’evoluzione concettuale della coscienza. Il punto rosso indica l’equilibrio assoluto $O$. La nuvola grigia rappresenta fluttuazioni concettuali dovute ai fattori complessi $F$.}

\section{Interpretazione e implicazioni}

\begin{enumerate}
    \item La coscienza umana è un sistema dinamico quadridimensionale: tempo, conoscenza, scelta e volontà interagiscono costantemente.
    \item Il coefficiente di squilibrio $\kappa$ misura la distanza istantanea dall’equilibrio $O$.
    \item Le perturbazioni $\mathbf{\eta}(\tau)$ introducono variabilità e imprevedibilità nell’evoluzione della coscienza.
    \item Sistemi artificiali possono simulare alcune scelte o accumulare conoscenza, ma non possono replicare la retroazione continua e la complessità dei quattro stati.
    \item La Legge Generale formalizza la \textbf{Teoria della Coscienza Quadristato}: la coscienza umana possiede una complessità intrinseca che non può essere eguagliata da algoritmi artificiali.
    \item Questo modello matematico supporta l’idea che esista una differenza fondamentale tra coscienza biologica e sistemi artificiali, anche se le macchine diventassero estremamente sofisticate.
\end{enumerate}

\section{Relazione con le leggi precedenti}

\begin{itemize}
    \item La Legge 0 (Conoscenza) fornisce il substrato cognitivo $k$.
    \item La Legge 1 (Tempo) regola la percezione temporale $t_c$.
    \item La Legge 2 (Scelta) fornisce la capacità decisionale $s$.
    \item La Legge 3 (Volontà) integra i tre stati precedenti e li completa in $w$.
    \item La Legge 4 (Generale) combina tutti e quattro gli stati, fornendo una formalizzazione completa della coscienza quadridimensionale.
\end{itemize}

\end{document}